\section{Checking Stateful Instances}
\label{sec:approach}
Figure~\ref{fig:overview} presents the workflow. The execution path constructs
states from the selected inputs and C bodies. The checking path turns those
states into immutable obligations, applies automatic reasoning, and submits
remaining goals to proof-assistant completion. Three interfaces connect the
paths: the entry domain, copied observations, and verdict evidence.

\subsection{Instantiate inputs and execute bodies}
Typed bindings specialize the root precondition. A binding can fix a logical
integer while leaving an address or part of a data structure symbolic. Entry
values remain immutable as C variables and memory change. The entry checker
establishes that the resulting domain contains at least one state, using
automatic reasoning or an existential Rocq obligation. Execution retains the
whole selected domain; an existence witness does not choose the sole test heap.

\tespec{} reuses \qcp{} to evaluate C expressions, generate arithmetic and
memory-safety obligations, and update the symbolic heap. Known guards select
a successor. Unresolved guards retain each successor whose infeasibility has
not been established. A call evaluates arguments, allocates fresh parameter
slots, and enters the available body. The callee receives fresh logical entry
bindings, including in recursive calls. On return, shared writes and the actual
return value survive; callee-local slots are retired. A loop repeatedly executes
its condition and body, including \texttt{break}, \texttt{continue}, and early
return. Resource exhaustion leaves the execution incomplete.

For example, a helper that increments a cell from 3 to 4 leaves the caller with
4 even if its contract promises only an existential integer greater than 3.
The summary is checked against that execution. It does not replace the
computed value. This is the distinction from simply strengthening the
precondition of a modular verification run with an input equality.

\subsection{Observe states without changing them}
\label{sec:boundaries}
At a selected boundary, let $P$ describe the reached state and $A$ the resolved
annotation. The checker emits
\[
             P\ \vdash\ A*F_{P,A}
\]
on a copy and continues execution with $P$. The frame $F_{P,A}$ contains
exactly the eligible local storage omitted by the annotation. Callee contracts
and visited loop invariants are observed in this way. A false assertion adds
a failing obligation; it does not prune the following execution.

\paragraph{Preserve the checked footprint.}
The memory model represents scalar C locals as owned cells. An auxiliary
annotation can omit unrelated local slots, so the checker resolves its
program variables and goal witnesses before selecting the local frame.
It preserves the omitted slots' types and current values. If a local pointer
holds $\&i$, resolution exposes the $i$ cell before framing and avoids counting
it twice. User-owned heap cells are never added merely because the assertion
omits them. In particular, a missing returned resource remains observable.

Let $m$ identify eligible local addresses, and split $h$ into $h_m$ and
$h_{\neg m}$. For extensional $A$, the scalar-cell model establishes
\[
 (A*\mathit{exact}(h_m))(h)\quad\Longleftrightarrow\quad A(h_{\neg m}).
\]
For an existential goal, the mask may depend on its witness. A successful
witness closes the existential obligation; failure for one chosen witness
does not refute all choices.

\paragraph{Select observation boundaries.}
The controller records which function, loop, and intermediate boundaries the
test checks. Both supported execution policies run available bodies and loops;
the richer policy also checks selected internal annotations. Named contract
alternatives are resolved by explicit routing, or retained as a disjunction
when no alternative is selected. An unresolved boundary match leaves a coverage
gap, which prevents complete acceptance.

\subsection{Joint heap and value normalization}
\label{sec:normalization-algorithm}
\label{sec:normalization}
The continuation can refer to logical values outside the current heap formula.
Eliminating $r$ from $\exists r.\ r=2\land\mathit{cell}(p,1)$ while leaving a
return term that refers to $r$ loses its value constraint. The normalization
interface therefore preserves the assertion $P$ together with every live term
$\bar t$, including results and saved call values. Define
\[
\obs(P,\bar t)(s,\bar v)=\exists\eta.\ P(s,\eta)\land
                  \bar v=\operatorname{eval}(\bar t)_{s,\eta}.
\]
Each admitted rewrite preserves this relation. Its composition consequently
preserves the observations used by entry and boundary checks.

\begin{lstlisting}[language={},caption={Normalization of the state and its live values.},label={lst:normalize}]
Normalize(Gamma, P, liveTerms, budget):
  while budget > 0 and an admitted reduction exists:
    resolve its definition and interpretation
    establish its guard; freshen typed local witnesses
    replace the occurrence by the complete equivalent body
    update dependent liveTerms consistently
    simplify exposed computations
    budget := budget - 1
  return (P, liveTerms)
\end{lstlisting}

\paragraph{Admitting reductions.}
Closed computations and known constructors expose reducible expressions.
Definition unfolding retains every pure conjunct, spatial resource, and typed
witness. For a guarded clause, activation requires
$g(\bar a)\Rightarrow(H(\bar a)\equiv\exists\bar w.\,B(\bar a,\bar w))$
in the selected interpretation. The executor establishes $g$ from the current
state. Unresolved guards retain the predicate application. Budgeted rewriting
terminates while preserving the meaning of residual terms.

For optional external rules, the controller checks the equation and its
association with the program's qualified definitions, constructors, spatial
notation, and layout. This prevents a same-named imported definition from
changing the meaning of an accepted equation. Imported predicates otherwise
use QCP's existing strategies. A checked pure consequence can be exposed by
$H\equiv H\land G$, retaining the original spatial predicate and its witnesses.
This path supports local proof progress without requiring a full layout
translation.

\paragraph{Preserving witnesses and symbolic addresses.}
Existing live variables remain protected during expansion. An unused witness
of nominal type $T$ requires inhabitance:
$(\exists w:T.\,P)\equiv(\mathit{inhabited}(T)\land P)$ when $w$ is absent
from $P$. Integer and list types supply canonical witnesses; an arbitrary
nominal type retains its existence obligation. Shared-base offset reductions
and closed list slices simplify operations with symbolic pointers. Their
ordinary C bounds and arithmetic safety obligations remain in the VC interface.


\subsection{Discharge immutable obligations}
\label{sec:evidence}
Automatic entailment and SMT reasoning first attempt the generated goals.
The optional \texttt{--spatial-smt} path encodes supported memory predicates
for the solver; unsupported formulas retain their original obligations.
For an entailment $P\vdash A$, a valid countermodel search asks whether
$P\land\neg A$ is satisfiable. An unsatisfiable answer discharges the entailment
only for a sound encoding. A satisfiable answer supplies a candidate violation,
whose memory interpretation and path reachability must also be established.

Residual goals remain Rocq propositions. An optional proof agent or a developer
fills proofs against frozen goal modules. The manifest fixes source hashes,
logical dependencies, target interpretation, goal identities, and proof
polarity. Complete proof replay checks every original obligation, including
entry existence and automatically discharged goals, audits transitive
assumptions, and performs an independent kernel check. All pending obligations
from one test execution form one proof-completion job.

\subsection{Aggregate execution and proof evidence}
A test passes when its entry domain is nonempty, execution and selected
observations are complete, and every required obligation is valid. A reachable
safety or assertion violation yields failure. A proof timeout, an unsupported
encoding, or unfinished execution yields an unresolved component. The
aggregation retains separate reasons for an empty entry domain, a program
violation, and an infrastructure error.

The implementation labels automatic solver evidence separately from checked
Rocq certificates. A negative proof of a VC establishes a failed test only
with a reachable violating state. These interfaces connect the workflow to
the formal guarantees below; the C frontend, target layout, and VC generator
must implement the stated transition and coverage conditions.
